%% Included in decks: slidemage-manual-copy
%%

%% DE translation: des-DE-LaTeX-unter-Linux-installieren
%%

\begin{standardframe}{Installing LaTeX in Linux}
	\small
	\begin{columns}[c,onlytextwidth]
		\begin{column}{0.86\linewidth}
			SlideMage compiles every deck with LuaLaTeX through latexmk, so it needs a
			TeX distribution. On Linux that is TeX Live, either from the distribution's
			package manager or from the upstream installer.
		\end{column}
		\begin{column}{0.1\linewidth}
			\raggedleft{\fontsize{34}{34}\selectfont\faLinux}
		\end{column}
	\end{columns}
	\par
	\begin{columns}[T,onlytextwidth]
		\begin{column}{0.56\linewidth}
			\shrinklists{1}
			\textbf{Package manager} (simplest)
			\begin{itemize}
				\item Debian/Ubuntu: \SMCode{sudo apt install texlive-full latexmk}
				\item Fedora: \SMCode{sudo dnf install texlive-scheme-full latexmk}
				\item Arch: \SMCode{sudo pacman -S texlive texlive-lang}
				\item Often a year or two behind the current TeX Live, which is fine for SlideMage
			\end{itemize}
		\end{column}
		\begin{column}{0.4\linewidth}
			\shrinklists{1}
			\textbf{Upstream TeX Live} (\SMCode{tug.org/texlive})
			\begin{itemize}
				\item \SMCode{install-tl}, scheme \emph{full} (about 8\,GB), no root needed in your home folder
				\item Add its \SMCode{bin} folder to \SMCode{PATH}
				\item Step by step: \href{https://tug.org/texlive/quickinstall.html}{quick install guide}
			\end{itemize}
			\medskip
			\textbf{Either way}
			\begin{itemize}
				\item Test in a terminal: \SMCode{lualatex --version}, \SMCode{latexmk -v}
			\end{itemize}
		\end{column}
	\end{columns}
\end{standardframe}
